\chapter{Certification Exam Mapping and Study Guide}
\index{certification!exam mapping}\index{Data Engineer Associate}\index{Data Engineer Professional}%
This appendix maps the Databricks Certified Data Engineer exams---Associate and
Professional---onto the book's chapters, then gives the study strategy and self-assessment
checklist that Chapter~24 operationalizes. Exam blueprints are revised periodically; always
confirm the current domain list and weightings in the official exam guides before booking
\cite{databricksCertification,databricksAssociateExamGuide,databricksProfessionalExamGuide}.

\section{The Two Exams at a Glance}
\index{certification!Associate versus Professional}%
\begin{itemize}
    \item \textbf{Data Engineer Associate.} Foundational, multiple-choice: the Databricks
    platform, Delta Lake, ETL with Spark SQL and PySpark, incremental ingestion, production
    pipelines (jobs, DLT), and governance with Unity Catalog. Reachable at the end of
    Part~III of this book (Week~12), with Parts~I--II as prerequisite fluency
    \cite{databricksAssociateExamGuide}.
    \item \textbf{Data Engineer Professional.} The same surface at production depth: scenario
    questions that reward failure-mode knowledge---checkpoint semantics, expectation
    behavior, concurrency, recovery, and performance trade-offs. Expect it to assume
    everything in Parts~IV--VI: streaming, delivery, security, and cost discipline
    \cite{databricksProfessionalExamGuide}.
\end{itemize}

\section{Domain-to-Chapter Mapping}
\index{certification!domain mapping}%
The book's curriculum was built around the exam domains; \cref{tab:appendixe-domains}
consolidates the mapping from Chapter~24 (its \cref{tab:ch24-domains}) with the remedial lab
for each domain. Score each domain from mock-exam evidence, not feeling.

\begin{table}[H]
\centering
\small
\begin{tabular}{@{}p{4.6cm}p{3.2cm}p{5.0cm}@{}}
\toprule
\textbf{Exam domain} & \textbf{Chapters} & \textbf{Rebuild these labs if weak} \\
\midrule
Databricks platform \& compute & 0, 1, 6 & Cluster policies; warehouse sizing lab \\
Delta Lake & 2, 5 & Enforcement/evolution/RESTORE drill; optimization benchmark \\
Medallion design & 3 & Quarantine and reconciliation lab \\
ETL with Spark \& DLT & 4, 9, 10 & Skew remediation; expectations poison drill \\
Incremental ingestion \& CDC & 11 & Auto Loader drift drill; SCD2 flow \\
Production pipelines \& orchestration & 12, 23 & Repair-run drill; bundle CI/CD lab \\
Streaming & 13--16 & Restart proof; watermark/dedup lab \\
Governance \& security & 8, 21, 22 & Masking + negative tests; share revocation \\
ML awareness & 17--20 & Registry alias flow; drift drill \\
\bottomrule
\end{tabular}
\caption{Exam domains mapped to chapters and remedial labs (extends Chapter~24's mapping).}
\label{tab:appendixe-domains}
\end{table}

\subsection{Where the Exams Probe}
The scenario questions concentrate on exactly the judgment calls the labs teach by breaking
things:
\begin{itemize}
    \item \textbf{Recovery semantics:} what \texttt{VACUUM} destroys, when \texttt{RESTORE}
    beats re-running, what a new checkpoint implies (Chapters~2, 11, 13).
    \item \textbf{Quality behavior:} when \texttt{expect\_or\_drop} versus
    \texttt{expect\_or\_fail} is correct, and what each does to downstream datasets
    (Chapter~10).
    \item \textbf{Concurrency and idempotency:} optimistic \texttt{MERGE} conflicts, safe
    repair runs, replay without duplicates (Chapters~2, 12, 14, 15).
    \item \textbf{Cost and layout:} job versus all-purpose clusters, liquid clustering versus
    partitioning, triggered versus continuous streaming (Chapters~1, 5, 13--16).
    \item \textbf{Governance mechanics:} grant scope, masking enforcement, lineage queries,
    sharing revocation (Chapters~8, 21, 22).
\end{itemize}

\section{Study Strategy}
\index{study strategy}\index{mock exam!strategy}%
\begin{enumerate}
    \item \textbf{Map confidence to domains.} Score every domain in
    \cref{tab:appendixe-domains} from evidence: the last lab you rebuilt unaided and your
    mock results, per domain.
    \item \textbf{Remediate by rebuilding, not re-reading.} A weak domain means rebuild that
    domain's lab from scratch, in a fresh schema, without notes. Reviewing answer keys
    without rebuilding converts a mock into trivia.
    \item \textbf{Take mocks under exam conditions.} Timed, closed-book, no notes, scored by
    domain. Classify every miss as knowledge gap (rebuild the lab), reasoning error
    (re-read two scenario questions slowly), or time management (budget per question).
    \item \textbf{Apply the two-pass standard.} Book the exam when you hold two consecutive
    80\%+ mocks with no domain below 70\%. One good mock is variance; two is signal.
    \item \textbf{Study failure modes, not feature lists.} For each domain, name its three
    most failure-mode-flavored facts before exam day.
\end{enumerate}

\section{Self-Assessment Checklist}
\index{self-assessment}%
Before booking, you should be able to answer ``yes, with evidence'' to each line:

\subsection*{Platform and Storage}
\begin{itemize}
    \item[$\square$] I can explain the control/data plane split and choose between all-purpose,
    job, and serverless compute for a given workload (Chapters~0, 1).
    \item[$\square$] I can walk a Delta commit end to end and resolve a \texttt{MERGE}
    conflict (Chapter~2).
    \item[$\square$] I have executed a corruption-and-\texttt{RESTORE} drill and can prove row
    counts before and after (Chapter~2).
    \item[$\square$] I can defend a layout choice---liquid clustering, partitioning, Z-ORDER---with
    a benchmark table (Chapter~5).
\end{itemize}

\subsection*{Pipelines and Processing}
\begin{itemize}
    \item[$\square$] I can read \texttt{explain} and the Spark UI and fix skew with two
    different mechanisms (Chapters~4, 9).
    \item[$\square$] I can build a DLT pipeline with warn/drop/fail expectations and query its
    event log (Chapter~10).
    \item[$\square$] I can ingest with Auto Loader through a mid-stream schema change and apply
    CDC with an SCD2 sequence guard (Chapter~11).
    \item[$\square$] I have run a Workflow repair and can state what makes it safe
    (Chapter~12).
\end{itemize}

\subsection*{Streaming}
\begin{itemize}
    \item[$\square$] I can argue exactly-once end to end for a pipeline I built
    (Chapters~13--15).
    \item[$\square$] I can size a watermark from a lateness histogram and defend the drop
    horizon (Chapter~14).
    \item[$\square$] I have reconciled a replayed Kafka window against the primary table
    (Chapter~15).
\end{itemize}

\subsection*{Governance, Delivery, and Readiness}
\begin{itemize}
    \item[$\square$] I can prove least privilege with a grant matrix and a negative test
    (Chapters~8, 21).
    \item[$\square$] I can trace a metric to its sources with lineage and audit queries
    (Chapter~22).
    \item[$\square$] I can deploy a bundle across dev/staging/prod with a CI pipeline and roll
    it back (Chapter~23).
    \item[$\square$] I hold two consecutive 80\%+ mocks with no domain below 70\%
    (Chapter~24).
    \item[$\square$] I can present my capstone in 15 minutes with one number per box and two
    trade-offs defended two levels deep (Chapter~24).
\end{itemize}

\begin{pitfall}
\textbf{Booking on the first good mock.} A single 85\% is variance; the exam fee buys a
second data point you could have gotten for free. Meet the two-pass standard, and keep the
lab-rebuild habit afterward---it is also how senior engineers stay honest.
\end{pitfall}
